% TOP_PROBLEM: {"id": 7000029, "title": "Geometry of Curves and Surfaces — Problem 8.2", "category_id": 6, "category": {"id": 6, "name": "geometry", "display_name": "Geometry", "description": "Euclidean and non-Euclidean geometry, geometric structures.", "slug": "geometry", "order_index": 6, "created_at": "2026-07-31T15:26:25.671Z"}, "status": "partially_solved"}
% FIRST_POSTED: null
% ENTRY_KIND: statement_only
% Imported from: batch18.tex; UnsolvedMath ID 7000029
\documentclass[11pt]{article}
\usepackage[margin=1in]{geometry}
\usepackage{amsmath,amssymb,mathtools}
\usepackage{fontspec,unicode-math}
\setmainfont{Latin Modern Roman}
\setsansfont{Latin Modern Sans}
\setmathfont{Latin Modern Math}
\usepackage{xurl,hyperref}
\hypersetup{colorlinks=true,urlcolor=blue,linkcolor=blue}
\newcommand{\sref}[2]{\hyperlink{source:#1:#2}{[S#2]}}
\newcommand{\cataloguescope}[1]{\paragraph{Recorded mathematical question.}#1}
\begin{document}
% UnsolvedMath ID 7000029; source catalogue code AMR-069-0029
\setcounter{section}{7000028}
\section{Geometry of Curves and Surfaces — Problem 8.2}\label{7000029}
{\small\sffamily UnsolvedMath ID 7000029 · Geometry}
\subsection{Definitions and mathematical statement}

\paragraph{Shared notation.}$\mathbb N=\{1,2,\ldots\}$, and $\mathbb Z,\mathbb Q,\mathbb R,\mathbb C$ denote the integers, rationals, reals and complex numbers. Any other conventions are specified locally.


Let $f:M\to\mathbb R^3$ be a smooth regular immersion of a surface. Locally choose a unit normal $n$; the shape operator is $A=-dn$, and its eigenvalues $k_1,k_2$ are the principal curvatures. An umbilic is a point where $k_1=k_2$. Away from umbilics the two eigendirections of $A$ give the principal curvature line fields. Let $p$ be an isolated umbilic. On a sufficiently small positively oriented loop around $p$, express either principal line direction in a continuous oriented local tangent frame by an angle $\theta$, taken modulo $\pi$. A continuous lift of $\theta$ around the loop defines the line-field index
\[\operatorname{ind}_p=\frac{\theta(2\pi)-\theta(0)}{2\pi}\in\tfrac12\mathbb Z.\]
The two perpendicular principal fields have the same index. The immersion is smooth at the umbilic itself; a point where it is only $C^1$ and not $C^2$ is outside this question.
\paragraph{Statement recorded in the supplied source.}
Does every isolated umbilic of every smooth immersed surface in $\mathbb R^3$ have principal-line-field index at most $1$? \sref{7000029}{2}
\subsection{Short English statement}
Can the principal curvature directions wind more than once around an isolated umbilic of a smooth surface?
\subsection{Sources}
\begin{description}
\item[\hypertarget{source:7000029:1}{S1}] ULAM.AI and the UnsolvedMath Contributors, \emph{UnsolvedMath: A Curated Collection of Open Mathematics Problems}, record 7000029 (AMR-069-0029). CC BY 4.0. \url{https://huggingface.co/datasets/ulamai/UnsolvedMath}.
\item[\hypertarget{source:7000029:2}{S2}] Naoya Ando, Toshifumi Fujiyama and Masaaki Umehara, \emph{$C^1$-umbilics with arbitrarily high indices}, arXiv:1508.03685v1 (2015), Introduction, pp. 1--2. \url{https://arxiv.org/pdf/1508.03685v1}.
\item[\hypertarget{source:7000029:3}{S3}] Mohammad Ghomi and Ralph Howard, \emph{Normal curvatures of asymptotically constant graphs and Carath\'eodory\textquotesingle s conjecture}, arXiv:1101.3031v2 (2011), Introduction, p. 3. \url{https://arxiv.org/pdf/1101.3031v2}.
\end{description}
\label{end:7000029}
\end{document}
